% GENERATED FILE. Edit canonical lessons, then run npm run generate:books.
% canonical-source-hash: fnv1a64:bb9b407fb4a07641
% canonical-lessons: PA-C236-manjna, PA-C236-hanghalna, PA-C236-uladdna, PA-C236-lota, PA-C236-katora

\chapter{Scrubbing the Pots}
\label{ch:pa-a2-scrubbing-the-pots}
\hlchaptermodality{236}

\begin{chapteropening}
\textbf{By the end of this chapter:} \emph{I can talk about washing the dishes.}

Complete the last lesson of chapter 236: I can talk about washing the dishes.
\end{chapteropening}

\section[\texorpdfstring{\pa{ਮਾਂਜਣਾ} (mānjṇā)}{ਮਾਂਜਣਾ (mānjṇā)}]{\pa{ਮਾਂਜਣਾ} (mānjṇā) — to scrub (pots)}

\label{lesson:PA-C236-manjna}



\begin{quote}
Before the new one: say the Punjabi for mud plaster, then the Punjabi for to plaster.
\end{quote}

\subsection*{You'll want to know: \pa{ਮਾਂਜਣਾ}}
\textbf{\pa{ਮਾਂਜਣਾ}} — \emph{mānjṇā} — \textquotedblleft{}to scrub (pots)\textquotedblright{}.

\begin{quote}
\pa{ਮੈਂ} \pa{ਰਾਤ} \pa{ਨੂੰ} \pa{ਭਾਂਡੇ} \pa{ਮਾਂਜਦੀ} \pa{ਹਾਂ।}
\end{quote}

\begin{quote}
\emph{maiṁ rāt nū̃ bhāṇḍe mānjdī hāṁ.}
\end{quote}

\begin{quote}
At night I scrub the pots.
\end{quote}

Past (I …): \pa{ਮੈਂ} \pa{ਭਾਂਡੇ} \pa{ਮਾਂਜੇ}, \emph{maiṁ bhāṇḍe mānje}.

\begin{cousinweb}
\textbf{\pa{ਮਾਂਜਣਾ}} (\emph{mānjṇā}) is inherited from Sanskrit \emph{mārjati}, \textquotedblleft{}wipes off, rubs clean\textquotedblright{}; CDIAL lists Punjabi \emph{mā̃jṇā}, \textquotedblleft{}to rub\textquotedblright{} (CDIAL 10080). The pot-scouring sense is CDIAL's Lahnda gloss: \emph{mā̃jhaṇ}, \textquotedblleft{}to clean out a vessel with sand or ashes\textquotedblright{}.
\end{cousinweb}

\begin{grammarlens}[title={What you've built}]
One more word for this chapter.
\end{grammarlens}

\subsection*{Your turn}
\begin{itemize}
\raggedright
  \item \emph{Say it:} \emph{mānjṇā}
  \item \emph{Say it:} \emph{mānjṇā}, once more
  \item \emph{Say it:} \emph{maiṁ rāt nū̃ bhāṇḍe mānjdī hāṁ.} — At night I scrub the pots.
  \item \emph{Recall:} say the Punjabi for mud plaster, then the Punjabi for to plaster, then say \emph{mānjṇā} again
\end{itemize}

\begin{tcolorbox}[breakable,colback=teal!4,colframe=teal!35!black,title={Before you move on}]
What does \pa{ਮਾਂਜਣਾ} mean? (\textquotedblleft{}To scrub (pots)\textquotedblright{}.) Say it once more.
\end{tcolorbox}
\section[\texorpdfstring{\pa{ਹੰਘਾਲਣਾ} (haṅghālṇā)}{ਹੰਘਾਲਣਾ (haṅghālṇā)}]{\pa{ਹੰਘਾਲਣਾ} (haṅghālṇā) — to rinse}

\label{lesson:PA-C236-hanghalna}



\begin{quote}
Before the new one: say the Punjabi for to plaster, then the Punjabi for to scrub.
\end{quote}

\subsection*{You'll want to know: \pa{ਹੰਘਾਲਣਾ}}
\textbf{\pa{ਹੰਘਾਲਣਾ}} — \emph{haṅghālṇā} — \textquotedblleft{}to rinse\textquotedblright{}.

\begin{quote}
\pa{ਮੈਂ} \pa{ਗਲਾਸ} \pa{ਪਾਣੀ} \pa{ਨਾਲ} \pa{ਹੰਘਾਲਦਾ} \pa{ਹਾਂ।}
\end{quote}

\begin{quote}
\emph{maiṁ galās pāṇī nāl haṅghāldā hāṁ.}
\end{quote}

\begin{quote}
I rinse the glass with water.
\end{quote}

Past (I …): \pa{ਮੈਂ} \pa{ਗਲਾਸ} \pa{ਹੰਘਾਲਿਆ}, \emph{maiṁ galās haṅghāliā}.

\begin{cousinweb}
\textbf{\pa{ਹੰਘਾਲਣਾ}} (\emph{haṅghālṇā}) is inherited from a reconstructed \emph{*āghālayati}, \textquotedblleft{}rinses\textquotedblright{}; CDIAL lists Punjabi \emph{hãghālṇā}, \textquotedblleft{}to rinse\textquotedblright{} (CDIAL 1058).
\end{cousinweb}

\begin{grammarlens}[title={What you've built}]
One more word for this chapter.
\end{grammarlens}

\subsection*{Your turn}
\begin{itemize}
\raggedright
  \item \emph{Say it:} \emph{haṅghālṇā}
  \item \emph{Say it:} \emph{haṅghālṇā}, once more
  \item \emph{Say it:} \emph{maiṁ galās pāṇī nāl haṅghāldā hāṁ.} — I rinse the glass with water.
  \item \emph{Recall:} say the Punjabi for to plaster, then the Punjabi for to scrub, then say \emph{haṅghālṇā} again
  \item \emph{Write it:} copy \textbf{\pa{ਹੰਘਾਲਣਾ}} once, checking the whole word
\end{itemize}

\begin{tcolorbox}[breakable,colback=teal!4,colframe=teal!35!black,title={Before you move on}]
What does \pa{ਹੰਘਾਲਣਾ} mean? (\textquotedblleft{}To rinse\textquotedblright{}.) Say it once more.
\end{tcolorbox}
\section[\texorpdfstring{\pa{ਉਲੱਦਣਾ} (uladdṇā)}{ਉਲੱਦਣਾ (uladdṇā)}]{\pa{ਉਲੱਦਣਾ} (uladdṇā) — to tip out; to empty}

\label{lesson:PA-C236-uladdna}



\begin{quote}
Before the new one: say the Punjabi for to scrub, then the Punjabi for to rinse.
\end{quote}

\subsection*{You'll want to know: \pa{ਉਲੱਦਣਾ}}
\textbf{\pa{ਉਲੱਦਣਾ}} — \emph{uladdṇā} — \textquotedblleft{}to tip out; to empty\textquotedblright{}.

\begin{quote}
\pa{ਮੈਂ} \pa{ਬਾਲਟੀ} \pa{ਦਾ} \pa{ਗੰਦਾ} \pa{ਪਾਣੀ} \pa{ਬਾਹਰ} \pa{ਉਲੱਦਦੀ} \pa{ਹਾਂ।}
\end{quote}

\begin{quote}
\emph{maiṁ bālṭī dā gandā pāṇī bāhar uladddī hāṁ.}
\end{quote}

\begin{quote}
I tip the dirty water from the bucket outside.
\end{quote}

Past (I …): \pa{ਮੈਂ} \pa{ਪਾਣੀ} \pa{ਉਲੱਦਿਆ}, \emph{maiṁ pāṇī uladdiā}.

\begin{cousinweb}
\textbf{\pa{ਉਲੱਦਣਾ}} (\emph{uladdṇā}) is inherited from a reconstructed \emph{*ullardayati}, \textquotedblleft{}overloads\textquotedblright{}; CDIAL lists Punjabi \emph{uladdṇā}, \textquotedblleft{}to overturn, upset\textquotedblright{} (CDIAL 2372).
\end{cousinweb}

\begin{grammarlens}[title={What you've built}]
One more word for this chapter.
\end{grammarlens}

\subsection*{Your turn}
\begin{itemize}
\raggedright
  \item \emph{Say it:} \emph{uladdṇā}
  \item \emph{Say it:} \emph{uladdṇā}, once more
  \item \emph{Say it:} \emph{maiṁ bālṭī dā gandā pāṇī bāhar uladddī hāṁ.} — I tip the dirty water from the bucket outside.
  \item \emph{Recall:} say the Punjabi for to scrub, then the Punjabi for to rinse, then say \emph{uladdṇā} again
  \item \emph{Write it:} copy \textbf{\pa{ਉਲੱਦਣਾ}} once, checking the whole word
\end{itemize}

\begin{tcolorbox}[breakable,colback=teal!4,colframe=teal!35!black,title={Before you move on}]
What does \pa{ਉਲੱਦਣਾ} mean? (\textquotedblleft{}To tip out; to empty\textquotedblright{}.) Say it once more.
\end{tcolorbox}
\section[\texorpdfstring{\pa{ਲੋਟਾ} (loṭā)}{ਲੋਟਾ (loṭā)}]{\pa{ਲੋਟਾ} (loṭā) — a small water pot}

\label{lesson:PA-C236-lota}



\begin{quote}
Before the new one: say the Punjabi for to rinse, then the Punjabi for to tip out.
\end{quote}

\subsection*{You'll want to know: \pa{ਲੋਟਾ}}
\textbf{\pa{ਲੋਟਾ}} — \emph{loṭā} — \textquotedblleft{}a small water pot\textquotedblright{}.

\begin{quote}
\pa{ਲੋਟੇ} \pa{ਵਿੱਚ} \pa{ਪਾਣੀ} \pa{ਹੈ।}
\end{quote}

\begin{quote}
\emph{loṭe vich pāṇī hai.}
\end{quote}

\begin{quote}
There is water in the pot.
\end{quote}

A masculine noun.

\begin{cousinweb}
\textbf{\pa{ਲੋਟਾ}} (\emph{loṭā}) is inherited from a reconstructed \emph{*lōṭṭa}, \textquotedblleft{}waterpot\textquotedblright{}; CDIAL lists Punjabi \emph{loṭṭā}, \textquotedblleft{}waterpot\textquotedblright{} (CDIAL 11133).
\end{cousinweb}

\begin{grammarlens}[title={What you've built}]
One more word for this chapter.
\end{grammarlens}

\subsection*{Your turn}
\begin{itemize}
\raggedright
  \item \emph{Say it:} \emph{loṭā}
  \item \emph{Say it:} \emph{loṭā}, once more
  \item \emph{Say it:} \emph{loṭe vich pāṇī hai.} — There is water in the pot.
  \item \emph{Recall:} say the Punjabi for to rinse, then the Punjabi for to tip out, then say \emph{loṭā} again
  \item \emph{Write it:} copy \textbf{\pa{ਲੋਟਾ}} once, checking the whole word
\end{itemize}

\begin{tcolorbox}[breakable,colback=teal!4,colframe=teal!35!black,title={Before you move on}]
What does \pa{ਲੋਟਾ} mean? (\textquotedblleft{}A small water pot\textquotedblright{}.) Say it once more.
\end{tcolorbox}
\section[\texorpdfstring{\pa{ਕਟੋਰਾ} (kaṭorā)}{ਕਟੋਰਾ (kaṭorā)}]{\pa{ਕਟੋਰਾ} (kaṭorā) — a deep bowl}

\label{lesson:PA-C236-katora}



\begin{quote}
Before the new one: say the Punjabi for to tip out, then the Punjabi for a small water pot.
\end{quote}

\subsection*{You'll want to know: \pa{ਕਟੋਰਾ}}
\textbf{\pa{ਕਟੋਰਾ}} — \emph{kaṭorā} — \textquotedblleft{}a deep bowl\textquotedblright{}.

\begin{quote}
\pa{ਕਟੋਰੇ} \pa{ਵਿੱਚ} \pa{ਦਾਲ} \pa{ਹੈ।}
\end{quote}

\begin{quote}
\emph{kaṭore vich dāl hai.}
\end{quote}

\begin{quote}
There is dal in the deep bowl.
\end{quote}

A masculine noun.

\begin{cousinweb}
\textbf{\pa{ਕਟੋਰਾ}} (\emph{kaṭorā}) is inherited from a reconstructed \emph{*kaṭṭōra}, \textquotedblleft{}cup\textquotedblright{}; CDIAL lists Punjabi \emph{kaṭorā}, \textquotedblleft{}small open-mouthed cup\textquotedblright{} (CDIAL 2648).
\end{cousinweb}

\begin{grammarlens}[title={What you've built}]
One more word for this chapter.
\end{grammarlens}

\subsection*{Your turn}
\begin{itemize}
\raggedright
  \item \emph{Say it:} \emph{kaṭorā}
  \item \emph{Say it:} \emph{kaṭorā}, once more
  \item \emph{Say it:} \emph{kaṭore vich dāl hai.} — There is dal in the deep bowl.
  \item \emph{Recall:} say the Punjabi for to tip out, then the Punjabi for a small water pot, then say \emph{kaṭorā} again
  \item \emph{Write it:} copy \textbf{\pa{ਕਟੋਰਾ}} once, checking the whole word
\end{itemize}

\begin{tcolorbox}[breakable,colback=teal!4,colframe=teal!35!black,title={Before you move on}]
What does \pa{ਕਟੋਰਾ} mean? (\textquotedblleft{}A deep bowl\textquotedblright{}.) Say it once more.
\end{tcolorbox}
